\Qhead{2.2}{Распространение сигнала в диспергирующей среде}{Распространение сигнала в диспергирующей среде. Простейшие физические модели диспергирующих сред. Волновой пакет в первом и втором приближении теории дисперсии. Фазовая и групповая скорости. Параболическое уравнение для огибающей. Расплывание и компрессия импульсов. Поле в средах с временной. Дисперсионные соотношения Крамерса--Кронига и принцип причинности.}
\Subfile{2.2.1}{sections/02_waves/q02/s01}
\Subfile{2.2.2}{sections/02_waves/q02/s02}
\Subfile{2.2.3}{sections/02_waves/q02/s03}

\recap{Первое приближение переносит огибающую с групповой скоростью; второе даёт её искажение и квадратичную спектральную фазу. Компрессия компенсирует начальную фазу. Причинная память среды связывает дисперсию с поглощением.}
\board{$k(\omega_0+\Omega)$ до второго порядка; уравнение огибающей; ширина гауссова пакета; причинный интеграл и одна формула Крамерса–Кронига.}
\oral{Почему большое или отрицательное $v_g$ не доказывает сверхсветовую передачу информации?}
\bait{Можно оставить знак дисперсии: он меняет направление частотной модуляции, хотя незачирпованный гауссов пакет расширяется при обоих знаках.}
\sources{Цысарь~\cite{tsysar1}, лекции~7--8, PDF-с.~34--45; Белокопытов~\cite{belokopytov1}, разделы о дисперсии. Причинные интегралы и нормировка гауссова решения проверены прямым преобразованием Фурье.}
